% Хавсралт B

\chapter{Юзкейсийн дэлгэрэнгүй тодорхойлолтууд}
\label{AppendixUsecaseSpecs}

\section{Үлдсэн юзкейсийн тодорхойлолтууд}
\label{sec:appendix-usecase-specs}

\paragraph{}
Энэ хавсралтад үндсэн бүлэгт зай хэмнэх зорилгоор оруулаагүй UC-004--UC-021 юзкейсүүдийн дэлгэрэнгүй тодорхойлолтыг байрлуулав.

\providecommand{\usecasecreatedby}{\shortname}
\providecommand{\usecasecreateddateearly}{2026-02-27}
\providecommand{\usecasecreateddatelate}{2026-03-12}

\subsection{Хэрэглэгчийн жагсаалт харах (UC-004)}

Хүснэгт~\ref{tab:uc-004}-д Хэрэглэгчийн жагсаалт харах юзкейсийн дэлгэрэнгүй тодорхойлолтыг үзүүлэв.

\begin{table}[H]
    \centering
    \renewcommand{\arraystretch}{1.3}
    \caption{UC-004 юзкейсийн тодорхойлолт}
    \label{tab:uc-004}
    \begin{tabular}{|p{3.8cm}|p{10.2cm}|}
        \hline
        \textbf{Юзкейс} & Хэрэглэгчийн жагсаалт харах \\ \hline
        \textbf{ID} & UC-004 \\ \hline
        \textbf{Үүсгэсэн} & \usecasecreatedby \\ \hline
        \textbf{Үүсгэсэн огноо} & \usecasecreateddateearly \\ \hline
        \textbf{Үндсэн тоглогч} & Админ \\ \hline
        \textbf{Нэмэлт тоглогч} & Байхгүй \\ \hline
        \textbf{Товч тайлбар} & Админ системд бүртгэлтэй бүх хэрэглэгчийн үндсэн мэдээлэл болон төлөвийг харна. \\ \hline
        \textbf{Өмнөх нөхцөл} &
        \begin{enumerate}
            \item Админ системд нэвтэрсэн байх.
        \end{enumerate}
        \\ \hline
        \textbf{Үндсэн урсгал} &
        \begin{enumerate}
            \item Админ хэрэглэгч удирдах хэсгийг нээнэ.
            \item Систем хэрэглэгчдийн жагсаалтыг ачаална.
            \item Админ хэрэглэгч бүрийн нэр, и-мэйл, үүрэг, хадгалалтын мэдээллийг шалгана.
        \end{enumerate}
        \\ \hline
        \textbf{Дараах нөхцөл} &
        \begin{enumerate}
            \item Хэрэглэгчийн удирдлагын мэдээлэл дэлгэц дээр харагдсан байна.
        \end{enumerate}
        \\ \hline
        \textbf{Альтернатив урсгал} &
        \begin{enumerate}
            \item Админ жагсаалтаас сонгосон хэрэглэгч рүү засварлах эсвэл устгах дараагийн урсгалд шилжинэ.
        \end{enumerate}
        \\ \hline
    \end{tabular}
\end{table}

\subsection{Хэрэглэгчийн мэдээлэл өөрчлөх (UC-005)}

Хүснэгт~\ref{tab:uc-005}-д Хэрэглэгчийн мэдээлэл өөрчлөх юзкейсийн дэлгэрэнгүй тодорхойлолтыг үзүүлэв.

\begin{table}[H]
    \centering
    \renewcommand{\arraystretch}{1.3}
    \caption{UC-005 юзкейсийн тодорхойлолт}
    \label{tab:uc-005}
    \begin{tabular}{|p{3.8cm}|p{10.2cm}|}
        \hline
        \textbf{Юзкейс} & Хэрэглэгчийн мэдээлэл өөрчлөх \\ \hline
        \textbf{ID} & UC-005 \\ \hline
        \textbf{Үүсгэсэн} & \usecasecreatedby \\ \hline
        \textbf{Үүсгэсэн огноо} & \usecasecreateddateearly \\ \hline
        \textbf{Үндсэн тоглогч} & Админ \\ \hline
        \textbf{Нэмэлт тоглогч} & Байхгүй \\ \hline
        \textbf{Товч тайлбар} & Админ хэрэглэгчийн үндсэн мэдээлэл болон хадгалалтын тохиргоог шинэчилнэ. \\ \hline
        \textbf{Өмнөх нөхцөл} &
        \begin{enumerate}
            \item Админ системд нэвтэрсэн байх.
            \item Өөрчлөх хэрэглэгч бүртгэлтэй байх.
        \end{enumerate}
        \\ \hline
        \textbf{Үндсэн урсгал} &
        \begin{enumerate}
            \item Админ жагсаалтаас хэрэглэгчээ сонгоно.
            \item Өөрчлөх мэдээллүүдийг шинэ утгаар оруулна.
            \item Систем өгөгдлийн хүчинтэй эсэхийг шалгана.
            \item Систем шинэ мэдээллийг хадгална.
        \end{enumerate}
        \\ \hline
        \textbf{Дараах нөхцөл} &
        \begin{enumerate}
            \item Хэрэглэгчийн шинэ мэдээлэл системд хадгалагдсан байна.
        \end{enumerate}
        \\ \hline
        \textbf{Альтернатив урсгал} &
        \begin{enumerate}
            \item Хэрэв зөвхөн хадгалалтын хязгаар өөрчлөгдөж байвал бусад талбар өмнөх утгаараа үлдэнэ.
            \item Хэрэв шинэчилсэн мэдээлэл давхцвал систем өөрчлөлтийг хадгалахгүй.
        \end{enumerate}
        \\ \hline
    \end{tabular}
\end{table}

\subsection{Хэрэглэгч устгах (UC-006)}

Хүснэгт~\ref{tab:uc-006}-д Хэрэглэгч устгах юзкейсийн дэлгэрэнгүй тодорхойлолтыг үзүүлэв.

\begin{table}[H]
    \centering
    \renewcommand{\arraystretch}{1.3}
    \caption{UC-006 юзкейсийн тодорхойлолт}
    \label{tab:uc-006}
    \begin{tabular}{|p{3.8cm}|p{10.2cm}|}
        \hline
        \textbf{Юзкейс} & Хэрэглэгч устгах \\ \hline
        \textbf{ID} & UC-006 \\ \hline
        \textbf{Үүсгэсэн} & \usecasecreatedby \\ \hline
        \textbf{Үүсгэсэн огноо} & \usecasecreateddateearly \\ \hline
        \textbf{Үндсэн тоглогч} & Админ \\ \hline
        \textbf{Нэмэлт тоглогч} & Байхгүй \\ \hline
        \textbf{Товч тайлбар} & Админ сонгосон энгийн хэрэглэгчийн бүртгэлийг системээс устгана. \\ \hline
        \textbf{Өмнөх нөхцөл} &
        \begin{enumerate}
            \item Админ системд нэвтэрсэн байх.
            \item Устгах хэрэглэгч админ биш байх.
        \end{enumerate}
        \\ \hline
        \textbf{Үндсэн урсгал} &
        \begin{enumerate}
            \item Админ жагсаалтаас хэрэглэгчээ сонгоно.
            \item Устгах үйлдлийг баталгаажуулна.
            \item Систем хэрэглэгчтэй холбоотой нэвтрэлтийн төлөв болон холбогдох өгөгдлүүдийг цэвэрлэнэ.
            \item Систем хэрэглэгчийн бүртгэлийг устгана.
        \end{enumerate}
        \\ \hline
        \textbf{Дараах нөхцөл} &
        \begin{enumerate}
            \item Устгасан хэрэглэгч системд нэвтрэх боломжгүй болно.
            \item Хэрэглэгч жагсаалтаас алга болсон байна.
        \end{enumerate}
        \\ \hline
        \textbf{Альтернатив урсгал} &
        \begin{enumerate}
            \item Хэрэв сонгосон хэрэглэгч олдохгүй бол систем үйлдлийг зогсооно.
            \item Хэрэв админ хэрэглэгчийг сонговол систем устгахыг зөвшөөрөхгүй.
        \end{enumerate}
        \\ \hline
    \end{tabular}
\end{table}

\subsection{Системд нэвтрэх (UC-007)}

Хүснэгт~\ref{tab:uc-007}-д Системд нэвтрэх юзкейсийн дэлгэрэнгүй тодорхойлолтыг үзүүлэв.

\begin{table}[H]
    \centering
    \renewcommand{\arraystretch}{1.3}
    \caption{UC-007 юзкейсийн тодорхойлолт}
    \label{tab:uc-007}
    \begin{tabular}{|p{3.8cm}|p{10.2cm}|}
        \hline
        \textbf{Юзкейс} & Системд нэвтрэх \\ \hline
        \textbf{ID} & UC-007 \\ \hline
        \textbf{Үүсгэсэн} & \usecasecreatedby \\ \hline
        \textbf{Үүсгэсэн огноо} & \usecasecreateddateearly \\ \hline
        \textbf{Үндсэн тоглогч} & Хэрэглэгч \\ \hline
        \textbf{Нэмэлт тоглогч} & Байхгүй \\ \hline
        \textbf{Товч тайлбар} & Хэрэглэгч нэвтрэх нэр, нууц үг ашиглан системд баталгаажиж үндсэн орчинд хандана. \\ \hline
        \textbf{Өмнөх нөхцөл} &
        \begin{enumerate}
            \item Хэрэглэгч системд бүртгэлтэй байх.
            \item Хэрэглэгчийн бүртгэл идэвхтэй байх.
        \end{enumerate}
        \\ \hline
        \textbf{Үндсэн урсгал} &
        \begin{enumerate}
            \item Хэрэглэгч нэвтрэх хуудсыг нээнэ.
            \item Нэвтрэх нэр болон нууц үгээ оруулна.
            \item Систем оруулсан мэдээллийг шалгана.
            \item Систем хэрэглэгчийг баталгаажуулж session үүсгэнэ.
            \item Хэрэглэгч өөрийн үндсэн хуудас руу шилжинэ.
        \end{enumerate}
        \\ \hline
        \textbf{Дараах нөхцөл} &
        \begin{enumerate}
            \item Идэвхтэй session үүссэн байна.
            \item Хэрэглэгч эрхтэй үйлдлүүддээ хандах боломжтой болно.
        \end{enumerate}
        \\ \hline
        \textbf{Альтернатив урсгал} &
        \begin{enumerate}
            \item Нууц үг буруу бол систем алдааны мэдээлэл үзүүлж дахин оролдох боломж олгоно.
            \item Админ хэрэглэгч амжилттай нэвтэрвэл админд зориулсан удирдлагын боломжууд мөн нээгдэнэ.
        \end{enumerate}
        \\ \hline
    \end{tabular}
\end{table}

\subsection{Нөөцүүдийг харах (UC-008)}

Хүснэгт~\ref{tab:uc-008}-д Нөөцүүдийг харах юзкейсийн дэлгэрэнгүй тодорхойлолтыг үзүүлэв.

\begin{table}[H]
    \centering
    \renewcommand{\arraystretch}{1.3}
    \caption{UC-008 юзкейсийн тодорхойлолт}
    \label{tab:uc-008}
    \begin{tabular}{|p{3.8cm}|p{10.2cm}|}
        \hline
        \textbf{Юзкейс} & Нөөцүүдийг харах \\ \hline
        \textbf{ID} & UC-008 \\ \hline
        \textbf{Үүсгэсэн} & \usecasecreatedby \\ \hline
        \textbf{Үүсгэсэн огноо} & \usecasecreateddateearly \\ \hline
        \textbf{Үндсэн тоглогч} & Хэрэглэгч \\ \hline
        \textbf{Нэмэлт тоглогч} & Байхгүй \\ \hline
        \textbf{Товч тайлбар} & Хэрэглэгч өөрийн болон өөртэй нь хуваалцсан файл, хавтсыг жагсаалтаар харна. \\ \hline
        \textbf{Өмнөх нөхцөл} &
        \begin{enumerate}
            \item Хэрэглэгч системд нэвтэрсэн байх.
        \end{enumerate}
        \\ \hline
        \textbf{Үндсэн урсгал} &
        \begin{enumerate}
            \item Хэрэглэгч нөөцийн үндсэн хэсгийг нээнэ.
            \item Систем тухайн байршил дахь файл, хавтасны жагсаалтыг харуулна.
            \item extension point: Нэрээр хайх (нэрээр хайх шаардлагатай бол)
            \item Хэрэглэгч хавтас сонгож дараагийн түвшний агуулга руу шилжинэ.
            \item Систем сонгосон байршлын агуулгыг шинэчилж харуулна.
        \end{enumerate}
        \\ \hline
        \textbf{Дараах нөхцөл} &
        \begin{enumerate}
            \item Сонгосон байршлын агуулга хэрэглэгчид харагдсан байна.
        \end{enumerate}
        \\ \hline
        \textbf{Альтернатив урсгал} &
        \begin{enumerate}
            \item Хэрэв хэрэглэгч нэрээр хайлт хийх шаардлагатай бол UC-009 ``Нэрээр хайх'' юзкейс ажиллана.
            \item Хэрэглэгч өөртэй нь хуваалцсан нөөцийн хэсэг рүү шилжин тухайн эрхийн хүрээнд агуулгыг харж болно.
        \end{enumerate}
        \\ \hline
    \end{tabular}
\end{table}

\subsection{Нэрээр хайх (UC-009)}

Хүснэгт~\ref{tab:uc-009}-д Нэрээр хайх юзкейсийн дэлгэрэнгүй тодорхойлолтыг үзүүлэв.

\begin{table}[H]
    \centering
    \renewcommand{\arraystretch}{1.3}
    \caption{UC-009 юзкейсийн тодорхойлолт}
    \label{tab:uc-009}
    \begin{tabular}{|p{3.8cm}|p{10.2cm}|}
        \hline
        \textbf{Өргөтгөх юзкейс} & Нэрээр хайх \\ \hline
        \textbf{ID} & UC-009 \\ \hline
        \textbf{Үүсгэсэн} & \usecasecreatedby \\ \hline
        \textbf{Үүсгэсэн огноо} & \usecasecreateddatelate \\ \hline
        \textbf{Үндсэн тоглогч} & Хэрэглэгч \\ \hline
        \textbf{Нэмэлт тоглогч} & Байхгүй \\ \hline
        \textbf{Товч тайлбар} & Хэрэглэгч өөрийн хандах эрхтэй нөөцүүдээс нэрийн дагуу хайлт хийнэ. \\ \hline
        \textbf{Сегмент 1 өмнөх нөхцөл} &
        \begin{enumerate}
            \item UC-008 ``Нөөцүүдийг харах'' юзкейсийн дэлгэц нээгдсэн байх.
            \item Нэрээр хайх шаардлага үүссэн байх.
        \end{enumerate}
        \\ \hline
        \textbf{Сегмент урсгал} &
        \begin{enumerate}
            \item Хэрэглэгч хайлтын талбарт түлхүүр үг оруулна.
            \item Систем нөөцийн нэрсээс тохирох үр дүнг хайна.
            \item Систем үр дүнг жагсаалтаар харуулна.
            \item Хэрэглэгч үр дүнгээс хүссэн нөөцөө сонгоно.
        \end{enumerate}
        \\ \hline
        \textbf{Сегмент 1 дараах нөхцөл} &
        \begin{enumerate}
            \item Хайлтын үр дүн хэрэглэгчид харагдсан байна.
        \end{enumerate}
        \\ \hline
        \textbf{Альтернатив урсгал} &
        \begin{enumerate}
            \item Хэрэв илэрц олдохгүй бол систем хоосон үр дүн үзүүлнэ.
            \item Хэрэв хайлтын талбар хоосорвол систем ердийн жагсаалт руу буцна.
        \end{enumerate}
        \\ \hline
    \end{tabular}
\end{table}

\subsection{Шинэ хавтас үүсгэх (UC-010)}

Хүснэгт~\ref{tab:uc-010}-д Шинэ хавтас үүсгэх юзкейсийн дэлгэрэнгүй тодорхойлолтыг үзүүлэв.

\begin{table}[H]
    \centering
    \renewcommand{\arraystretch}{1.3}
    \caption{UC-010 юзкейсийн тодорхойлолт}
    \label{tab:uc-010}
    \begin{tabular}{|p{3.8cm}|p{10.2cm}|}
        \hline
        \textbf{Юзкейс} & Шинэ хавтас үүсгэх \\ \hline
        \textbf{ID} & UC-010 \\ \hline
        \textbf{Үүсгэсэн} & \usecasecreatedby \\ \hline
        \textbf{Үүсгэсэн огноо} & \usecasecreateddateearly \\ \hline
        \textbf{Үндсэн тоглогч} & Хэрэглэгч \\ \hline
        \textbf{Нэмэлт тоглогч} & Байхгүй \\ \hline
        \textbf{Товч тайлбар} & Хэрэглэгч зөвшөөрөгдсөн байршилд шинэ хавтас үүсгэнэ. \\ \hline
        \textbf{Өмнөх нөхцөл} &
        \begin{enumerate}
            \item Хэрэглэгч системд нэвтэрсэн байх.
            \item Зорилтот байршил дээр үүсгэх эрхтэй байх.
        \end{enumerate}
        \\ \hline
        \textbf{Үндсэн урсгал} &
        \begin{enumerate}
            \item Хэрэглэгч хавтас үүсгэх байршлаа сонгоно.
            \item Шинэ хавтас үүсгэх үйлдлийг дарна.
            \item Шинэ хавтасны нэрийг оруулна.
            \item Систем нэрийн хүчинтэй байдлыг шалгана.
            \item Систем хавтсыг үүсгэж жагсаалтад нэмнэ.
        \end{enumerate}
        \\ \hline
        \textbf{Дараах нөхцөл} &
        \begin{enumerate}
            \item Шинэ хавтас сонгосон байршилд үүссэн байна.
        \end{enumerate}
        \\ \hline
        \textbf{Альтернатив урсгал} &
        \begin{enumerate}
            \item Хэрэв ижил нэртэй хавтас байгаа бол систем дахин нэр оруулахыг шаардана.
            \item Хэрэв зорилтот байршил root түвшин бол хавтас хэрэглэгчийн үндсэн орчинд үүснэ.
        \end{enumerate}
        \\ \hline
    \end{tabular}
\end{table}

\subsection{Файл байршуулах (UC-011)}

Хүснэгт~\ref{tab:uc-011}-д Файл байршуулах юзкейсийн дэлгэрэнгүй тодорхойлолтыг үзүүлэв.

\begin{table}[H]
    \centering
    \renewcommand{\arraystretch}{1.3}
    \caption{UC-011 юзкейсийн тодорхойлолт}
    \label{tab:uc-011}
    \begin{tabular}{|p{3.8cm}|p{10.2cm}|}
        \hline
        \textbf{Юзкейс} & Файл байршуулах \\ \hline
        \textbf{ID} & UC-011 \\ \hline
        \textbf{Үүсгэсэн} & \usecasecreatedby \\ \hline
        \textbf{Үүсгэсэн огноо} & \usecasecreateddateearly \\ \hline
        \textbf{Үндсэн тоглогч} & Хэрэглэгч \\ \hline
        \textbf{Нэмэлт тоглогч} & Байхгүй \\ \hline
        \textbf{Товч тайлбар} & Хэрэглэгч зөвшөөрөгдсөн хавтас дотор файл байршуулж хадгална. \\ \hline
        \textbf{Өмнөх нөхцөл} &
        \begin{enumerate}
            \item Хэрэглэгч системд нэвтэрсэн байх.
            \item Зорилтот хавтас дээр байршуулах эрхтэй байх.
        \end{enumerate}
        \\ \hline
        \textbf{Үндсэн урсгал} &
        \begin{enumerate}
            \item Хэрэглэгч файл байршуулах хавтсаа сонгоно.
            \item Файлаа сонгож байршуулах үйлдлийг эхлүүлнэ.
            \item Систем эрх, багтаамж, файлын хэмжээ болон нэрийг шалгана.
            \item Систем файлыг хадгална.
            \item Систем жагсаалтыг шинэчилж шинээр байршуулсан файлыг харуулна.
        \end{enumerate}
        \\ \hline
        \textbf{Дараах нөхцөл} &
        \begin{enumerate}
            \item Файл сонгосон байршилд хадгалагдсан байна.
            \item Хадгалалтын ашиглалтын хэмжээ шинэчлэгдэнэ.
        \end{enumerate}
        \\ \hline
        \textbf{Альтернатив урсгал} &
        \begin{enumerate}
            \item Хэрэв багтаамж хүрэлцэхгүй бол систем байршуулахыг зогсооно.
            \item Хэрэв дамжуулалт тасарвал хэрэглэгч дахин байршуулах боломжтой байна.
        \end{enumerate}
        \\ \hline
    \end{tabular}
\end{table}

\subsection{Нөөцийн нэр өөрчлөх (UC-012)}

Хүснэгт~\ref{tab:uc-012}-д Нөөцийн нэр өөрчлөх юзкейсийн дэлгэрэнгүй тодорхойлолтыг үзүүлэв.

\begin{table}[H]
    \centering
    \renewcommand{\arraystretch}{1.3}
    \caption{UC-012 юзкейсийн тодорхойлолт}
    \label{tab:uc-012}
    \begin{tabular}{|p{3.8cm}|p{10.2cm}|}
        \hline
        \textbf{Юзкейс} & Нөөцийн нэр өөрчлөх \\ \hline
        \textbf{ID} & UC-012 \\ \hline
        \textbf{Үүсгэсэн} & \usecasecreatedby \\ \hline
        \textbf{Үүсгэсэн огноо} & \usecasecreateddateearly \\ \hline
        \textbf{Үндсэн тоглогч} & Хэрэглэгч \\ \hline
        \textbf{Нэмэлт тоглогч} & Байхгүй \\ \hline
        \textbf{Товч тайлбар} & Хэрэглэгч өөрийн эрхтэй файл эсвэл хавтсын нэрийг шинэчилнэ. \\ \hline
        \textbf{Өмнөх нөхцөл} &
        \begin{enumerate}
            \item Хэрэглэгч системд нэвтэрсэн байх.
            \item Сонгосон нөөц дээр нэр өөрчлөх эрхтэй байх.
        \end{enumerate}
        \\ \hline
        \textbf{Үндсэн урсгал} &
        \begin{enumerate}
            \item Хэрэглэгч нэр өөрчлөх нөөцөө сонгоно.
            \item Шинэ нэрийг оруулна.
            \item Систем нэрийн дүрэм болон давхардлыг шалгана.
            \item Систем шинэ нэрийг хадгална.
        \end{enumerate}
        \\ \hline
        \textbf{Дараах нөхцөл} &
        \begin{enumerate}
            \item Нөөц шинэ нэрээрээ жагсаалтад харагдана.
        \end{enumerate}
        \\ \hline
        \textbf{Альтернатив урсгал} &
        \begin{enumerate}
            \item Хэрэв шинэ нэр хоосон эсвэл хүчингүй бол систем өөрчлөлтийг хадгалахгүй.
            \item Хэрэв ижил нэртэй нөөц байвал систем дахин нэр сонгохыг шаардана.
        \end{enumerate}
        \\ \hline
    \end{tabular}
\end{table}

\subsection{Нөөц сэргээх (UC-013)}

Хүснэгт~\ref{tab:uc-013}-д Нөөц сэргээх юзкейсийн дэлгэрэнгүй тодорхойлолтыг үзүүлэв.

\begin{table}[H]
    \centering
    \renewcommand{\arraystretch}{1.3}
    \caption{UC-013 юзкейсийн тодорхойлолт}
    \label{tab:uc-013}
    \begin{tabular}{|p{3.8cm}|p{10.2cm}|}
        \hline
        \textbf{Юзкейс} & Нөөц сэргээх \\ \hline
        \textbf{ID} & UC-013 \\ \hline
        \textbf{Үүсгэсэн} & \usecasecreatedby \\ \hline
        \textbf{Үүсгэсэн огноо} & \usecasecreateddatelate \\ \hline
        \textbf{Үндсэн тоглогч} & Хэрэглэгч \\ \hline
        \textbf{Нэмэлт тоглогч} & Байхгүй \\ \hline
        \textbf{Товч тайлбар} & Эзэмшигч хэрэглэгч устгасан төлөвт байгаа файл эсвэл хавтсаа сэргээнэ. \\ \hline
        \textbf{Өмнөх нөхцөл} &
        \begin{enumerate}
            \item Хэрэглэгч системд нэвтэрсэн байх.
            \item Сэргээх нөөц нь устгасан төлөвтэй байх.
        \end{enumerate}
        \\ \hline
        \textbf{Үндсэн урсгал} &
        \begin{enumerate}
            \item Хэрэглэгч устгасан нөөцийн хэсгийг нээнэ.
            \item Сэргээх файл эсвэл хавтсаа сонгоно.
            \item Сэргээх үйлдлийг баталгаажуулна.
            \item Систем нөөцийг өмнөх идэвхтэй төлөвт буцаана.
        \end{enumerate}
        \\ \hline
        \textbf{Дараах нөхцөл} &
        \begin{enumerate}
            \item Нөөц идэвхтэй жагсаалт руу буцаж орсон байна.
        \end{enumerate}
        \\ \hline
        \textbf{Альтернатив урсгал} &
        \begin{enumerate}
            \item Хэрэв өмнөх байршилд ижил нэртэй нөөц үүссэн бол систем шинэ нэр сонгох эсвэл өөр байршил заахыг шаардаж болно.
            \item Хэрэв хэрэглэгч сэргээхээс татгалзвал үйлдэл цуцлагдана.
        \end{enumerate}
        \\ \hline
    \end{tabular}
\end{table}

\subsection{Нөөц шилжүүлэх (UC-014)}

Хүснэгт~\ref{tab:uc-014}-д Нөөц шилжүүлэх юзкейсийн дэлгэрэнгүй тодорхойлолтыг үзүүлэв.

\begin{table}[H]
    \centering
    \renewcommand{\arraystretch}{1.3}
    \caption{UC-014 юзкейсийн тодорхойлолт}
    \label{tab:uc-014}
    \begin{tabular}{|p{3.8cm}|p{10.2cm}|}
        \hline
        \textbf{Юзкейс} & Нөөц шилжүүлэх \\ \hline
        \textbf{ID} & UC-014 \\ \hline
        \textbf{Үүсгэсэн} & \usecasecreatedby \\ \hline
        \textbf{Үүсгэсэн огноо} & \usecasecreateddateearly \\ \hline
        \textbf{Үндсэн тоглогч} & Хэрэглэгч \\ \hline
        \textbf{Нэмэлт тоглогч} & Байхгүй \\ \hline
        \textbf{Товч тайлбар} & Эзэмшигч хэрэглэгч файл эсвэл хавтсаа өөр байршил руу шилжүүлнэ. \\ \hline
        \textbf{Өмнөх нөхцөл} &
        \begin{enumerate}
            \item Хэрэглэгч системд нэвтэрсэн байх.
            \item Шилжүүлэх нөөц болон зорилтот хавтас дээр шаардлагатай эрхтэй байх.
        \end{enumerate}
        \\ \hline
        \textbf{Үндсэн урсгал} &
        \begin{enumerate}
            \item Хэрэглэгч нэг буюу хэд хэдэн нөөцөө сонгоно.
            \item Шилжүүлэх үйлдлийг эхлүүлнэ.
            \item Зорилтот хавтсыг сонгоно.
            \item Систем бүтэц болон эрхийн шалгалт хийнэ.
            \item Систем нөөцийг шинэ байршилд шилжүүлнэ.
        \end{enumerate}
        \\ \hline
        \textbf{Дараах нөхцөл} &
        \begin{enumerate}
            \item Нөөц шинэ байршилд харагдана.
            \item Хуучин байршлаас сонгосон нөөц алга болсон байна.
        \end{enumerate}
        \\ \hline
        \textbf{Альтернатив урсгал} &
        \begin{enumerate}
            \item Хэрэв нэг дор олон нөөц сонгосон бол систем бүгдийг хамтад нь шилжүүлнэ.
            \item Хэрэв хавтсыг өөрийнх нь доторх дэд хавтас руу шилжүүлэх оролдлого гарвал систем үйлдлийг зогсооно.
        \end{enumerate}
        \\ \hline
    \end{tabular}
\end{table}

\subsection{Нөөц устгах (UC-015)}

Хүснэгт~\ref{tab:uc-015}-д Нөөц устгах юзкейсийн дэлгэрэнгүй тодорхойлолтыг үзүүлэв.

\begin{table}[H]
    \centering
    \renewcommand{\arraystretch}{1.3}
    \caption{UC-015 юзкейсийн тодорхойлолт}
    \label{tab:uc-015}
    \begin{tabular}{|p{3.8cm}|p{10.2cm}|}
        \hline
        \textbf{Юзкейс} & Нөөц устгах \\ \hline
        \textbf{ID} & UC-015 \\ \hline
        \textbf{Үүсгэсэн} & \usecasecreatedby \\ \hline
        \textbf{Үүсгэсэн огноо} & \usecasecreateddateearly \\ \hline
        \textbf{Үндсэн тоглогч} & Хэрэглэгч \\ \hline
        \textbf{Нэмэлт тоглогч} & Байхгүй \\ \hline
        \textbf{Товч тайлбар} & Эзэмшигч хэрэглэгч өөрийн файл эсвэл хавтсыг идэвхтэй жагсаалтаас устгасан төлөв рүү шилжүүлнэ. \\ \hline
        \textbf{Өмнөх нөхцөл} &
        \begin{enumerate}
            \item Хэрэглэгч системд нэвтэрсэн байх.
            \item Сонгосон нөөц тухайн хэрэглэгчийн эзэмшлийнх байх.
        \end{enumerate}
        \\ \hline
        \textbf{Үндсэн урсгал} &
        \begin{enumerate}
            \item Хэрэглэгч устгах нөөцөө сонгоно.
            \item Устгах үйлдлийг баталгаажуулна.
            \item Систем нөөцийг идэвхтэй жагсаалтаас хасаж устгасан төлөвт шилжүүлнэ.
            \item extension point: Нөөцийг бүр мөсөн устгах (сэргээх боломжгүйгээр устгахыг сонговол)
        \end{enumerate}
        \\ \hline
        \textbf{Дараах нөхцөл} &
        \begin{enumerate}
            \item Нөөц устгасан хэсэгт хадгалагдана.
            \item Идэвхтэй жагсаалтад харагдахгүй болно.
        \end{enumerate}
        \\ \hline
        \textbf{Альтернатив урсгал} &
        \begin{enumerate}
            \item Хэрэв хэрэглэгч бүр мөсөн устгах шаардлагатай бол UC-016 ``Нөөцийг бүр мөсөн устгах'' юзкейс рүү шилжинэ.
            \item Хэрэглэгч нэг дор олон нөөц сонгон устгаж болно.
        \end{enumerate}
        \\ \hline
    \end{tabular}
\end{table}

\subsection{Нөөцийг бүр мөсөн устгах (UC-016)}

Хүснэгт~\ref{tab:uc-016}-д Нөөцийг бүр мөсөн устгах юзкейсийн дэлгэрэнгүй тодорхойлолтыг үзүүлэв.

\begin{table}[H]
    \centering
    \renewcommand{\arraystretch}{1.3}
    \caption{UC-016 юзкейсийн тодорхойлолт}
    \label{tab:uc-016}
    \begin{tabular}{|p{3.8cm}|p{10.2cm}|}
        \hline
        \textbf{Өргөтгөх юзкейс} & Нөөцийг бүр мөсөн устгах \\ \hline
        \textbf{ID} & UC-016 \\ \hline
        \textbf{Үүсгэсэн} & \usecasecreatedby \\ \hline
        \textbf{Үүсгэсэн огноо} & \usecasecreateddatelate \\ \hline
        \textbf{Үндсэн тоглогч} & Хэрэглэгч \\ \hline
        \textbf{Нэмэлт тоглогч} & Байхгүй \\ \hline
        \textbf{Товч тайлбар} & Устгасан төлөвт байгаа нөөцийг сэргээх боломжгүйгээр системээс бүрэн устгана. \\ \hline
        \textbf{Сегмент 1 өмнөх нөхцөл} &
        \begin{enumerate}
            \item UC-015 ``Нөөц устгах'' юзкейсийн extension point дээр ирсэн байх.
            \item Нөөц устгасан төлөвт байгаа байх.
            \item Хэрэглэгч сэргээх боломжгүйгээр устгахыг сонгосон байх.
        \end{enumerate}
        \\ \hline
        \textbf{Сегмент урсгал} &
        \begin{enumerate}
            \item Хэрэглэгч устгасан нөөцийн хэсгээс нөөцөө сонгоно.
            \item Бүр мөсөн устгах үйлдлийг сонгоно.
            \item Систем сэрэмжлүүлэх мэдээлэл харуулна.
            \item Хэрэглэгч баталгаажуулна.
            \item Систем нөөцийг физик түвшинд устгана.
        \end{enumerate}
        \\ \hline
        \textbf{Сегмент 1 дараах нөхцөл} &
        \begin{enumerate}
            \item Нөөц системээс бүрэн арилсан байна.
            \item Тухайн нөөцийг сэргээх боломжгүй болно.
        \end{enumerate}
        \\ \hline
        \textbf{Альтернатив урсгал} &
        \begin{enumerate}
            \item Хэрэв хэрэглэгч баталгаажуулахгүй бол үйлдэл цуцлагдана.
            \item Хэрэв нөөц олдохгүй бол систем алдааны мэдээлэл харуулна.
        \end{enumerate}
        \\ \hline
    \end{tabular}
\end{table}

\subsection{Файл татаж авах (UC-017)}

Хүснэгт~\ref{tab:uc-017}-д Файл татаж авах юзкейсийн дэлгэрэнгүй тодорхойлолтыг үзүүлэв.

\begin{table}[H]
    \centering
    \renewcommand{\arraystretch}{1.3}
    \caption{UC-017 юзкейсийн тодорхойлолт}
    \label{tab:uc-017}
    \begin{tabular}{|p{3.8cm}|p{10.2cm}|}
        \hline
        \textbf{Юзкейс} & Файл татаж авах \\ \hline
        \textbf{ID} & UC-017 \\ \hline
        \textbf{Үүсгэсэн} & \usecasecreatedby \\ \hline
        \textbf{Үүсгэсэн огноо} & \usecasecreateddateearly \\ \hline
        \textbf{Үндсэн тоглогч} & Хэрэглэгч \\ \hline
        \textbf{Нэмэлт тоглогч} & Байхгүй \\ \hline
        \textbf{Товч тайлбар} & Хэрэглэгч өөрийн хандах эрхтэй файлыг төхөөрөмж дээрээ татаж авна. \\ \hline
        \textbf{Өмнөх нөхцөл} &
        \begin{enumerate}
            \item Хэрэглэгч системд нэвтэрсэн байх.
            \item Сонгосон файлд хандах эрхтэй байх.
        \end{enumerate}
        \\ \hline
        \textbf{Үндсэн урсгал} &
        \begin{enumerate}
            \item Хэрэглэгч татах файлаа сонгоно.
            \item Татаж авах үйлдлийг эхлүүлнэ.
            \item Систем хандалтын эрх болон файлын бэлэн байдлыг шалгана.
            \item extension point: Шахсан файл хэлбэрээр татах (нэгээс олон файл эсвэл хавтас сонгосон бол)
            \item Систем файлыг хэрэглэгч рүү дамжуулна.
        \end{enumerate}
        \\ \hline
        \textbf{Дараах нөхцөл} &
        \begin{enumerate}
            \item Файл хэрэглэгчийн төхөөрөмжид татагдсан байна.
        \end{enumerate}
        \\ \hline
        \textbf{Альтернатив урсгал} &
        \begin{enumerate}
            \item Хэрэв хэд хэдэн файл эсвэл хавтас хамт татах шаардлагатай бол UC-018 ``Шахсан файл хэлбэрээр татах'' юзкейс ашиглагдана.
            \item Хэрэв дамжуулалт тасарвал хэрэглэгч дахин татах боломжтой байна.
        \end{enumerate}
        \\ \hline
    \end{tabular}
\end{table}

\subsection{Шахсан файл хэлбэрээр татах (UC-018)}

Хүснэгт~\ref{tab:uc-018}-д Шахсан файл хэлбэрээр татах юзкейсийн дэлгэрэнгүй тодорхойлолтыг үзүүлэв.

\begin{table}[H]
    \centering
    \renewcommand{\arraystretch}{1.3}
    \caption{UC-018 юзкейсийн тодорхойлолт}
    \label{tab:uc-018}
    \begin{tabular}{|p{3.8cm}|p{10.2cm}|}
        \hline
        \textbf{Өргөтгөх юзкейс} & Шахсан файл хэлбэрээр татах \\ \hline
        \textbf{ID} & UC-018 \\ \hline
        \textbf{Үүсгэсэн} & \usecasecreatedby \\ \hline
        \textbf{Үүсгэсэн огноо} & \usecasecreateddatelate \\ \hline
        \textbf{Үндсэн тоглогч} & Хэрэглэгч \\ \hline
        \textbf{Нэмэлт тоглогч} & Байхгүй \\ \hline
        \textbf{Товч тайлбар} & Хэрэглэгч хэд хэдэн файл эсвэл хавтсыг нэг шахсан файл болгон татаж авна. \\ \hline
        \textbf{Сегмент 1 өмнөх нөхцөл} &
        \begin{enumerate}
            \item UC-017 ``Файл татаж авах'' юзкейсийн extension point дээр ирсэн байх.
            \item Нэгээс олон файл эсвэл хавтас сонгогдсон байх.
            \item Сонгосон бүх нөөцөд хандах эрхтэй байх.
        \end{enumerate}
        \\ \hline
        \textbf{Сегмент урсгал} &
        \begin{enumerate}
            \item Хэрэглэгч нэгээс олон файл эсвэл хавтас сонгоно.
            \item Шахсан хэлбэрээр татах үйлдлийг сонгоно.
            \item Систем сонгосон нөөцүүдийг түр архив файл болгон бэлтгэнэ.
            \item Систем бэлтгэсэн архивыг хэрэглэгч рүү татахаар дамжуулна.
        \end{enumerate}
        \\ \hline
        \textbf{Сегмент 1 дараах нөхцөл} &
        \begin{enumerate}
            \item Хэрэглэгч нэг шахсан файл хүлээн авсан байна.
        \end{enumerate}
        \\ \hline
        \textbf{Альтернатив урсгал} &
        \begin{enumerate}
            \item Хэрэв сонголтод зөвхөн нэг файл байвал ердийн таталтын урсгал ашиглаж болно.
            \item Хэрэв шахалтын явцад алдаа гарвал систем архив үүсгэхгүй бөгөөд хэрэглэгчид мэдээлнэ.
        \end{enumerate}
        \\ \hline
    \end{tabular}
\end{table}

\subsection{Нөөц хуваалцах (UC-019)}

Хүснэгт~\ref{tab:uc-019}-д Нөөц хуваалцах юзкейсийн дэлгэрэнгүй тодорхойлолтыг үзүүлэв.

\begin{table}[H]
    \centering
    \renewcommand{\arraystretch}{1.3}
    \caption{UC-019 юзкейсийн тодорхойлолт}
    \label{tab:uc-019}
    \begin{tabular}{|p{3.8cm}|p{10.2cm}|}
        \hline
        \textbf{Юзкейс} & Нөөц хуваалцах \\ \hline
        \textbf{ID} & UC-019 \\ \hline
        \textbf{Үүсгэсэн} & \usecasecreatedby \\ \hline
        \textbf{Үүсгэсэн огноо} & \usecasecreateddateearly \\ \hline
        \textbf{Үндсэн тоглогч} & Хэрэглэгч \\ \hline
        \textbf{Нэмэлт тоглогч} & Байхгүй \\ \hline
        \textbf{Товч тайлбар} & Эзэмшигч хэрэглэгч өөрийн файл эсвэл хавтсыг бусад хэрэглэгчидтэй хуваалцана. \\ \hline
        \textbf{Өмнөх нөхцөл} &
        \begin{enumerate}
            \item Хэрэглэгч системд нэвтэрсэн байх.
            \item Хуваалцах нөөц нь тухайн хэрэглэгчийн эзэмшлийнх байх.
        \end{enumerate}
        \\ \hline
        \textbf{Үндсэн урсгал} &
        \begin{enumerate}
            \item Хэрэглэгч хуваалцах нөөцөө сонгоно.
            \item Хуваалцах үйлдлийг сонгоно.
            \item Систем зорилтот хэрэглэгч сонгох боломж харуулна.
            \item include (Хандалтын түвшин сонгох)
            \item Систем хуваалцах мэдээллийг хадгална.
        \end{enumerate}
        \\ \hline
        \textbf{Дараах нөхцөл} &
        \begin{enumerate}
            \item Сонгосон хэрэглэгч тухайн нөөцөд оноосон түвшний эрхтэй болно.
        \end{enumerate}
        \\ \hline
        \textbf{Альтернатив урсгал} &
        \begin{enumerate}
            \item Хэрэв өмнө нь эрх олгосон хэрэглэгчийг дахин сонговол систем өмнөх эрхийн түвшинг шинэчилнэ.
            \item Хэрэв зорилтот хэрэглэгч олдохгүй бол үйлдэл цуцлагдана.
        \end{enumerate}
        \\ \hline
    \end{tabular}
\end{table}

\subsection{Хандалтын түвшин сонгох (UC-020)}

Хүснэгт~\ref{tab:uc-020}-д Хандалтын түвшин сонгох юзкейсийн дэлгэрэнгүй тодорхойлолтыг үзүүлэв.

\begin{table}[H]
    \centering
    \renewcommand{\arraystretch}{1.3}
    \caption{UC-020 юзкейсийн тодорхойлолт}
    \label{tab:uc-020}
    \begin{tabular}{|p{3.8cm}|p{10.2cm}|}
        \hline
        \textbf{Юзкейс} & Хандалтын түвшин сонгох \\ \hline
        \textbf{ID} & UC-020 \\ \hline
        \textbf{Үүсгэсэн} & \usecasecreatedby \\ \hline
        \textbf{Үүсгэсэн огноо} & \usecasecreateddatelate \\ \hline
        \textbf{Үндсэн тоглогч} & Хэрэглэгч \\ \hline
        \textbf{Нэмэлт тоглогч} & Байхгүй \\ \hline
        \textbf{Товч тайлбар} & Нөөц хуваалцах үед зорилтот хэрэглэгчид олгох эрхийн түвшинг сонгоно. \\ \hline
        \textbf{Өмнөх нөхцөл} &
        \begin{enumerate}
            \item UC-019 ``Нөөц хуваалцах'' юзкейс эхэлсэн байх.
            \item Хуваалцах хэрэглэгч сонгогдсон байх.
        \end{enumerate}
        \\ \hline
        \textbf{Үндсэн урсгал} &
        \begin{enumerate}
            \item Хэрэглэгч харах эсвэл засварлах эрхийн нэгийг сонгоно.
            \item Систем сонгосон түвшинг харуулж баталгаажуулна.
            \item Систем тухайн түвшний эрхийг хуваалцах мэдээлэлд хадгална.
        \end{enumerate}
        \\ \hline
        \textbf{Дараах нөхцөл} &
        \begin{enumerate}
            \item Зорилтот хэрэглэгчид тодорхой хандалтын түвшин оноогдсон байна.
        \end{enumerate}
        \\ \hline
        \textbf{Альтернатив урсгал} &
        \begin{enumerate}
            \item Хэрэв хэрэглэгч эрхийн түвшнээ өөрчлөхөөр шийдвэл баталгаажуулахын өмнө дахин сонголт хийж болно.
            \item Хэрэв сонголт хийгдээгүй бол хуваалцах үйлдэл дуусахгүй.
        \end{enumerate}
        \\ \hline
    \end{tabular}
\end{table}

\subsection{Онцлох төлөв өөрчлөх (UC-021)}

Хүснэгт~\ref{tab:uc-021}-д Онцлох төлөв өөрчлөх юзкейсийн дэлгэрэнгүй тодорхойлолтыг үзүүлэв.

\begin{table}[H]
    \centering
    \renewcommand{\arraystretch}{1.3}
    \caption{UC-021 юзкейсийн тодорхойлолт}
    \label{tab:uc-021}
    \begin{tabular}{|p{3.8cm}|p{10.2cm}|}
        \hline
        \textbf{Юзкейс} & Онцлох төлөв өөрчлөх \\ \hline
        \textbf{ID} & UC-021 \\ \hline
        \textbf{Үүсгэсэн} & \usecasecreatedby \\ \hline
        \textbf{Үүсгэсэн огноо} & \usecasecreateddatelate \\ \hline
        \textbf{Үндсэн тоглогч} & Хэрэглэгч \\ \hline
        \textbf{Нэмэлт тоглогч} & Байхгүй \\ \hline
        \textbf{Товч тайлбар} & Хэрэглэгч өөрийн нөөцийг онцолсон төлөвт оруулах эсвэл тухайн төлөвөөс гаргана. \\ \hline
        \textbf{Өмнөх нөхцөл} &
        \begin{enumerate}
            \item Хэрэглэгч системд нэвтэрсэн байх.
            \item Сонгосон нөөц нь тухайн хэрэглэгчийн эзэмшлийнх байх.
        \end{enumerate}
        \\ \hline
        \textbf{Үндсэн урсгал} &
        \begin{enumerate}
            \item Хэрэглэгч файл эсвэл хавтсаа сонгоно.
            \item Онцлох төлөв өөрчлөх үйлдлийг хийнэ.
            \item Систем тухайн нөөцийн онцолсон төлөвийг сольж хадгална.
        \end{enumerate}
        \\ \hline
        \textbf{Дараах нөхцөл} &
        \begin{enumerate}
            \item Нөөцийн онцолсон төлөв шинэчлэгдсэн байна.
        \end{enumerate}
        \\ \hline
        \textbf{Альтернатив урсгал} &
        \begin{enumerate}
            \item Хэрэв нөөц өмнө нь онцолсон байсан бол систем онцолсон төлөвөөс гаргана.
            \item Хэрэв нөөц онцлоогүй байсан бол систем онцолсон жагсаалтад нэмнэ.
        \end{enumerate}
        \\ \hline
    \end{tabular}
\end{table}
